% GENERATED FILE. Edit canonical lessons, then run npm run generate:books.
% canonical-source-hash: fnv1a64:5749966aad412684
% canonical-lessons: KA-C359-utsava, KA-C359-meravanige, KA-C359-ghosavakya, KA-C359-meresu, KA-C359-dhvaja

\chapter{The Town Festival}
\label{ch:ka-the-town-festival}
\hlchaptermodality{359}

\begin{chapteropening}
\textbf{By the end of this chapter:} \emph{I can say a public festival, a procession, a slogan, to parade and a flag.}

Complete the last lesson of chapter 359: I can say a public festival, a procession, a slogan, to parade and a flag.
\end{chapteropening}

\section[\texorpdfstring{\kn{ಉತ್ಸವ} (utsava)}{ಉತ್ಸವ (utsava)}]{\kn{ಉತ್ಸವ} (utsava) — a public festival; a celebration}

\label{lesson:KA-C359-utsava}



\begin{quote}
Before the new one: say the Kannada for bail, then the Kannada for to confess.
\end{quote}

\subsection*{You'll want to know: \kn{ಉತ್ಸವ}}
\textbf{\kn{ಉತ್ಸವ}} — \emph{utsava} — \textquotedblleft{}a public festival; a celebration\textquotedblright{}.

\begin{quote}
\kn{ಊರಿನ} \kn{ಉತ್ಸವಕ್ಕೆ} \kn{ಹಳ್ಳಿಗಳಿಂದ} \kn{ಜನ} \kn{ಬಂದರು}.
\end{quote}

\begin{quote}
\emph{ūrina utsavakke haḷḷigaḷinda jana bandaru.}
\end{quote}

\begin{quote}
People came from the villages to the town festival.
\end{quote}

Plural: \kn{ಉತ್ಸವಗಳು}, \emph{utsavagaḷu}.

A public festival, often of a temple or a town; \kn{ಹಬ್ಬ} is the family festival.

\begin{cousinweb}
\textbf{\kn{ಉತ್ಸವ}} (\emph{utsava}) is borrowed from Sanskrit \emph{utsava}, \textquotedblleft{}an enterprise\textquotedblright{}, and later \textquotedblleft{}a festival\textquotedblright{} (CDIAL 1876).
\end{cousinweb}

\begin{grammarlens}[title={What you've built}]
One more word for telling what happened.
\end{grammarlens}

\subsection*{Your turn}
\begin{itemize}
\raggedright
  \item \emph{Say it:} \emph{utsava}
  \item \emph{Say it:} \emph{utsava}, once more
  \item \emph{Say it:} \emph{ūrina utsavakke haḷḷigaḷinda jana bandaru.} — People came from the villages to the town festival.
  \item \emph{Recall:} say the Kannada for bail, then the Kannada for to confess, then say \emph{utsava} again
\end{itemize}

\begin{tcolorbox}[breakable,colback=teal!4,colframe=teal!35!black,title={Before you move on}]
What does \kn{ಉತ್ಸವ} mean? (\textquotedblleft{}A public festival; a celebration\textquotedblright{}.) Say it once more.
\end{tcolorbox}
\section[\texorpdfstring{\kn{ಮೆರವಣಿಗೆ} (meravaṇige)}{ಮೆರವಣಿಗೆ (meravaṇige)}]{\kn{ಮೆರವಣಿಗೆ} (meravaṇige) — a procession}

\label{lesson:KA-C359-meravanige}



\begin{quote}
Before the new one: say the Kannada for to confess, then the Kannada for a public festival.
\end{quote}

\subsection*{You'll want to know: \kn{ಮೆರವಣಿಗೆ}}
\textbf{\kn{ಮೆರವಣಿಗೆ}} — \emph{meravaṇige} — \textquotedblleft{}a procession\textquotedblright{}.

\begin{quote}
\kn{ಮೆರವಣಿಗೆ} \kn{ಮುಖ್ಯ} \kn{ರಸ್ತೆಯಲ್ಲಿ} \kn{ನಿಧಾನವಾಗಿ} \kn{ಸಾಗಿತು}.
\end{quote}

\begin{quote}
\emph{meravaṇige mukhya rasteyalli nidhānavāgi sāgitu.}
\end{quote}

\begin{quote}
The procession moved slowly along the main road.
\end{quote}

Plural: \kn{ಮೆರವಣಿಗೆಗಳು}, \emph{meravaṇigegaḷu}.

A procession \textquotedblleft{}moves on\textquotedblright{}: \kn{ಮೆರವಣಿಗೆ} \kn{ಸಾಗಿತು}.

\begin{cousinweb}
\textbf{\kn{ಮೆರವಣಿಗೆ}} (\emph{meravaṇige}) is native Dravidian: DEDR 5074 lists Tamil \emph{meravaṇai}, \textquotedblleft{}a procession\textquotedblright{}, with Kannada \emph{merugu}, \textquotedblleft{}to shine\textquotedblright{}, and \emph{mere}, \textquotedblleft{}to parade\textquotedblright{}.
\end{cousinweb}

\begin{grammarlens}[title={What you've built}]
One more word for telling what happened.
\end{grammarlens}

\subsection*{Your turn}
\begin{itemize}
\raggedright
  \item \emph{Say it:} \emph{meravaṇige}
  \item \emph{Say it:} \emph{meravaṇige}, once more
  \item \emph{Say it:} \emph{meravaṇige mukhya rasteyalli nidhānavāgi sāgitu.} — The procession moved slowly along the main road.
  \item \emph{Recall:} say the Kannada for to confess, then the Kannada for a public festival, then say \emph{meravaṇige} again
\end{itemize}

\begin{tcolorbox}[breakable,colback=teal!4,colframe=teal!35!black,title={Before you move on}]
What does \kn{ಮೆರವಣಿಗೆ} mean? (\textquotedblleft{}A procession\textquotedblright{}.) Say it once more.
\end{tcolorbox}
\section[\texorpdfstring{\kn{ಘೋಷವಾಕ್ಯ} (ghōṣavākya)}{ಘೋಷವಾಕ್ಯ (ghōṣavākya)}]{\kn{ಘೋಷವಾಕ್ಯ} (ghōṣavākya) — a slogan}

\label{lesson:KA-C359-ghosavakya}



\begin{quote}
Before the new one: say the Kannada for a public festival, then the Kannada for a procession.
\end{quote}

\subsection*{You'll want to know: \kn{ಘೋಷವಾಕ್ಯ}}
\textbf{\kn{ಘೋಷವಾಕ್ಯ}} — \emph{ghōṣavākya} — \textquotedblleft{}a slogan\textquotedblright{}.

\begin{quote}
\kn{ಮೆರವಣಿಗೆಯಲ್ಲಿ} \kn{ಜನ} \kn{ಘೋಷವಾಕ್ಯ} \kn{ಕೂಗಿದರು}.
\end{quote}

\begin{quote}
\emph{meravaṇigeyalli jana ghōṣavākya kūgidaru.}
\end{quote}

\begin{quote}
People shouted slogans in the procession.
\end{quote}

Plural: \kn{ಘೋಷವಾಕ್ಯಗಳು}, \emph{ghōṣavākyagaḷu}.

\kn{ಘೋಷ}, a cry, and \kn{ವಾಕ್ಯ}, a sentence.

\begin{cousinweb}
\textbf{\kn{ಘೋಷವಾಕ್ಯ}} (\emph{ghōṣavākya}) is borrowed from Sanskrit, a modern compound of \emph{ghoṣa}, \textquotedblleft{}a noise, a cry\textquotedblright{} (Monier-Williams), and \emph{vākya}, \textquotedblleft{}speech, words\textquotedblright{} (CDIAL 11468).
\end{cousinweb}

\begin{grammarlens}[title={What you've built}]
One more word for telling what happened.
\end{grammarlens}

\subsection*{Your turn}
\begin{itemize}
\raggedright
  \item \emph{Say it:} \emph{ghōṣavākya}
  \item \emph{Say it:} \emph{ghōṣavākya}, once more
  \item \emph{Say it:} \emph{meravaṇigeyalli jana ghōṣavākya kūgidaru.} — People shouted slogans in the procession.
  \item \emph{Recall:} say the Kannada for a public festival, then the Kannada for a procession, then say \emph{ghōṣavākya} again
\end{itemize}

\begin{tcolorbox}[breakable,colback=teal!4,colframe=teal!35!black,title={Before you move on}]
What does \kn{ಘೋಷವಾಕ್ಯ} mean? (\textquotedblleft{}A slogan\textquotedblright{}.) Say it once more.
\end{tcolorbox}
\section[\texorpdfstring{\kn{ಮೆರೆಸು} (meresu)}{ಮೆರೆಸು (meresu)}]{\kn{ಮೆರೆಸು} (meresu) — to parade; to carry in procession}

\label{lesson:KA-C359-meresu}



\begin{quote}
Before the new one: say the Kannada for a procession, then the Kannada for a slogan.
\end{quote}

\subsection*{You'll want to know: \kn{ಮೆರೆಸು}}
\textbf{\kn{ಮೆರೆಸು}} — \emph{meresu} — \textquotedblleft{}to parade; to carry in procession\textquotedblright{}.

\begin{quote}
\kn{ಜನರು} \kn{ಗೆದ್ದ} \kn{ಆಟಗಾರನನ್ನು} \kn{ಬೀದಿಯಲ್ಲಿ} \kn{ಮೆರೆಸಿದರು}.
\end{quote}

\begin{quote}
\emph{janaru gedda āṭagāranannu bīdiyalli meresidaru.}
\end{quote}

\begin{quote}
The people paraded the winning player through the streets.
\end{quote}

Past (I …): \kn{ಮೆರೆಸಿದೆ}, \emph{mereside}.

Someone honoured is carried or led through the streets with music.

\begin{cousinweb}
\textbf{\kn{ಮೆರೆಸು}} (\emph{meresu}) is native Dravidian: the causative of \emph{mere}, \textquotedblleft{}to shine, to show off\textquotedblright{}, which DEDR 5074 lists with Kannada \emph{merugu}, \textquotedblleft{}to shine\textquotedblright{}.
\end{cousinweb}

\begin{grammarlens}[title={What you've built}]
One more word for telling what happened.
\end{grammarlens}

\subsection*{Your turn}
\begin{itemize}
\raggedright
  \item \emph{Say it:} \emph{meresu}
  \item \emph{Say it:} \emph{meresu}, once more
  \item \emph{Say it:} \emph{janaru gedda āṭagāranannu bīdiyalli meresidaru.} — The people paraded the winning player through the streets.
  \item \emph{Recall:} say the Kannada for a procession, then the Kannada for a slogan, then say \emph{meresu} again
\end{itemize}

\begin{tcolorbox}[breakable,colback=teal!4,colframe=teal!35!black,title={Before you move on}]
What does \kn{ಮೆರೆಸು} mean? (\textquotedblleft{}To parade; to carry in procession\textquotedblright{}.) Say it once more.
\end{tcolorbox}
\section[\texorpdfstring{\kn{ಧ್ವಜ} (dhvaja)}{ಧ್ವಜ (dhvaja)}]{\kn{ಧ್ವಜ} (dhvaja) — a flag; a banner}

\label{lesson:KA-C359-dhvaja}



\begin{quote}
Before the new one: say the Kannada for a slogan, then the Kannada for to parade.
\end{quote}

\subsection*{You'll want to know: \kn{ಧ್ವಜ}}
\textbf{\kn{ಧ್ವಜ}} — \emph{dhvaja} — \textquotedblleft{}a flag; a banner\textquotedblright{}.

\begin{quote}
\kn{ಉತ್ಸವದ} \kn{ದಿನ} \kn{ಎಲ್ಲ} \kn{ಮನೆಗಳ} \kn{ಮೇಲೆ} \kn{ಧ್ವಜ} \kn{ಹಾರುತ್ತಿತ್ತು}.
\end{quote}

\begin{quote}
\emph{utsavada dina ella manegaḷa mēle dhvaja hāruttittu.}
\end{quote}

\begin{quote}
On the festival day a flag flew on every house.
\end{quote}

Plural: \kn{ಧ್ವಜಗಳು}, \emph{dhvajagaḷu}.

A flag is \textquotedblleft{}hoisted\textquotedblright{}: \kn{ಧ್ವಜ} \kn{ಹಾರಿಸು}. \kn{ರಾಷ್ಟ್ರಧ್ವಜ} is the national flag.

\begin{cousinweb}
\textbf{\kn{ಧ್ವಜ}} (\emph{dhvaja}) is borrowed from Sanskrit \emph{dhvaja}, \textquotedblleft{}a banner, a flag, a standard\textquotedblright{} (Monier-Williams).
\end{cousinweb}

\begin{grammarlens}[title={What you've built}]
One more word for telling what happened.
\end{grammarlens}

\subsection*{Your turn}
\begin{itemize}
\raggedright
  \item \emph{Say it:} \emph{dhvaja}
  \item \emph{Say it:} \emph{dhvaja}, once more
  \item \emph{Say it:} \emph{utsavada dina ella manegaḷa mēle dhvaja hāruttittu.} — On the festival day a flag flew on every house.
  \item \emph{Recall:} say the Kannada for a slogan, then the Kannada for to parade, then say \emph{dhvaja} again
\end{itemize}

\begin{tcolorbox}[breakable,colback=teal!4,colframe=teal!35!black,title={Before you move on}]
What does \kn{ಧ್ವಜ} mean? (\textquotedblleft{}A flag; a banner\textquotedblright{}.) Say it once more.
\end{tcolorbox}
